% Solutions of 20-backtracking.
\section*{Chapter \ref{ch:b1:backtracking} --- Exhaustive Search and Backtracking}

\begin{solution}{exo:b1:backtracking:1}
\code{subsets_binary} gives the subsets with \code{d} for the numbers 8 to 15, in this order: d, ad, bd, abd, cd, acd,
bcd, abcd. The tree decides \code{d} last and includes before it excludes, so the subsets with \code{d} are every other
\omterm{def:b1:induction-and-termination:rooted-tree}{leaf}, starting with the first: abcd, abd, acd, ad, bcd, bd, cd, d.
\end{solution}

\begin{solution}{exo:b1:backtracking:2}
The longest decreasing suffix is $5, 4, 2$, so $i = 1$ ($a[1] = 3$). The smallest larger element of the suffix is 4,
at $j = 3$. The exchange gives $[1, 4, 5, 3, 2]$, and reversing the suffix gives $[1, 4, 2, 3, 5]$.
\end{solution}

\begin{solution}{exo:b1:backtracking:3}
Level by level: 1 (the root), 4, 6, 4 and 2: 17 nodes, as in \cref{out:b1:backtracking:queens}. Four leaves are dead
ends: \code{0 2}, \code{0 3 1}, \code{3 0 2} and \code{3 1}.
\end{solution}

\begin{solution}{exo:b1:backtracking:4}
The remaining sums are $18, 11, 6, 2, 0$. The root and the node ``7 in'' pass both tests; the node ``7 in, 5 in'' has
total $12 > 9$ and is the first one pruned, by the test ``too much already''.
\end{solution}

\begin{solution}{exo:b1:backtracking:5}
Without the undo, the loop at position $k$ swaps with an \omterm{def:b1:pointers-and-arrays:array}{array} the recursive call has changed. After the first branch,
the \omterm{def:b1:pointers-and-arrays:array}{array} is \code{acb}, not \code{abc}; the swap for $i = 1$ brings \code{c} to the front instead of \code{b}, and
after that branch the \omterm{def:b1:pointers-and-arrays:array}{array} is \code{cba}, so the swap for $i = 2$ brings \code{a} back. The letter \code{b} never
reaches the front: \code{bac} and \code{bca} are missing, \code{abc} and \code{acb} come twice. The number of outputs is
still $3!$, so a test that only counts them misses the bug.
\end{solution}

\begin{solution}{exo:b1:backtracking:6}
Both solutions pass through a partial sum above 6 (8, and $5 + 3$) that a later negative value brings back: the
\omterm{def:b1:backtracking:pruning}{pruning} discards \omterm{def:b1:induction-and-termination:rooted-tree}{subtrees} that hold solutions. A sound test uses the extreme sums still reachable: with $N_i$ the sum
of the negative values among the undecided ones and $P_i$ that of the positive ones, prune when $\mathit{total} + N_i >
\mathit{target}$ or $\mathit{total} + P_i < \mathit{target}$. With positive values, $N_i = 0$ and this is the chapter's
test.
\end{solution}

\begin{solution}{exo:b1:backtracking:7}
\omcode{code/bachelor-1/20-backtracking/exercises/distinct_parts.py}{4}{15}{Partitions into distinct parts}
Parts are chosen in decreasing order, which makes each partition appear once. Parts at most \code{part} sum to at most
$\mathit{part}(\mathit{part} + 1)/2$; when that is below $n$, no smaller first part can do better, and the loop stops.
The test compares the counts for $n < 40$ with a dynamic-programming count (2 for $n = 3$, 10 for $n = 10$).
\end{solution}

\begin{solution}{exo:b1:backtracking:8}
\omcode{code/bachelor-1/20-backtracking/exercises/wordsearch.py}{5}{29}{A word-search solver}
The \omterm{def:b1:backtracking:backtracking}{partial solution} is the path; a choice is a neighbour of its last cell; the \omterm{def:b1:backtracking:pruning}{constraints} are the next letter and
the unused cell. The removal of the cell after a dead end is the undo.
\end{solution}

\begin{solution}{exo:b1:backtracking:9}
The nodes at depth $k$ are the sequences of $k$ distinct elements, $n(n - 1) \cdots (n - k + 1) = n!/(n - k)!$ of them.
Summing over $k$ and putting $j = n - k$ gives $\sum_{j=0}^{n} n!/j!$, which is less than $n! \sum_{j \geq 0} 1/j! = e
\cdot n!$. For $n = 4$: $24 + 24 + 12 + 4 + 1 = 65$.
\end{solution}

\begin{solution}{exo:b1:backtracking:10}
\omcode{code/bachelor-1/20-backtracking/exercises/colouring.py}{4}{26}{Three colours by \omterm{def:b1:backtracking:backtracking}{backtracking}}
The \omterm{def:b1:backtracking:backtracking}{partial solution} assigns colours to the first vertices of the order; a colour used by a coloured neighbour is never
tried. The test compares with the $3^8$ assignments of 20 random graphs, and colours the town's friendship graph (39
people, 50 friendships), which has no 2-colouring. The encoding of \cref{ch:b1:propositional-logic} turns the same
question into clauses for a general solver: nothing in it is specific to graphs, which is its strength, and nothing
uses the order of the vertices, which is the strength of the direct search.
\end{solution}

\begin{solution}{exo:b1:backtracking:11}
Let the remaining items be sorted by decreasing $v/w$, with room $R$. A completion chooses a set $S$ of them of weight at
most $R$; with $x_i = 1$ on $S$ and 0 elsewhere, it is a solution of the relaxed problem, so its value is at most the
relaxed optimum. The greedy fill is that optimum: in any relaxed solution with room left over or with weight on an item
of lower ratio while an item of higher ratio is not full, moving weight to the higher ratio does not lower the value; the
greedy fill is the one where no such move is left. So $\mathit{value} + \mathit{relaxation} \geq$ the value of every
completion, and a node with $\mathit{value} + \mathit{relaxation} \leq \mathit{best}$ has no completion better than the
best known: \omterm{def:b1:backtracking:pruning}{pruning} it loses no improvement, and the optimum is still found.
\end{solution}

\begin{solution}{exo:b1:backtracking:12}
The root has six children, the first queen in columns 0 to 5. The mirror $c \mapsto 5 - c$ maps the search below column
$c$ onto the search below column $5 - c$, node for node, so the \omterm{def:b1:induction-and-termination:rooted-tree}{subtrees} come in three equal pairs: $1 + 2(x + y + z)$.
\code{exercises/queens_subtrees.py} gives $x, y, z = 22, 29, 25$, and $1 + 2 \cdot 76 = 153$.
\end{solution}

\begin{solution}{pb:b1:backtracking:1}
\textbf{1.} Two queens on a row or a column attack each other, so each row holds at most one queen; eight queens on eight
rows then hold exactly one each. The same for columns.

\textbf{2.} With $p(r)$ the column of row $r$'s queen, distinct columns make $p$ a permutation. Two queens share a
diagonal when $r - c$ or $r + c$ agree, so these values must all be distinct.

\textbf{3.} $\binom{64}{8} = 4\,426\,165\,368$ placements; $8! = 40\,320$ with one queen per row and column.

\textbf{4.} $r - c$ runs from $-7$ to 7, $r + c$ from 0 to 14: 15 diagonals in each direction.

\textbf{5.} The test asks whether a value is in a set, which is constant time on average for a hash set
(\cref{ch:b1:hash-tables}): three such questions per candidate.

\textbf{6.} $8!$ is about 40 thousand permutations, each tested in $O(n)$; $14!$ is about $8.7 \times 10^{10}$, far
beyond what a Python loop goes through in reasonable time.

\textbf{7.} 2057 nodes with \omterm{def:b1:backtracking:pruning}{pruning}, 109\,601 without: 53 times fewer.

\textbf{8.} Nothing forces the counts to grow with $n$: on the $6 \times 6$ board the \omterm{def:b1:backtracking:pruning}{constraints} happen to leave only
4 solutions, the images of a single one. The counts grow on the whole, but not at every step.

\textbf{9.} It roughly doubles with each $n$ (3.8, 6.0, 12.8, 24.8, 53.3, 117.5, 277.6): the gain itself grows
exponentially, while both counts still grow exponentially.

\textbf{10.} No. Placing column by column on a board is placing row by row on the board turned about its diagonal, a
square board with the same \omterm{def:b1:backtracking:pruning}{constraints}: the tree is the same.

\textbf{11.} No queen is placed above the first row, so nothing attacks any of its cells; each later queen loses the
columns and the diagonals of the queens above it.

\textbf{12.} The node count grows by a factor that rises from about 3 to 4.2; with factors near 4.5, the estimate for
$n = 12$ is about $35\,539 \times 4.5^2 \approx 720\,000$. (The program gives 856\,189.)

\textbf{13.} With $m = 7$: the identity $(r, c)$; the quarter-turns $(c, m - r)$ and $(m - c, r)$; the half-turn $(m - r,
m - c)$; the mirrors in the two axes $(r, m - c)$ and $(m - r, c)$; the mirrors in the two diagonals $(c, r)$ and $(m -
c, m - r)$.

\textbf{14.} Each is a bijection of the cells that maps rows and columns to rows or columns, and diagonals of either
direction to diagonals. Two cells that share no line are mapped to two cells that share none: a solution goes to a
solution.

\textbf{15.} The eight symmetries are closed under composition and inverse. If $q = g(p)$, then the images of $q$ are
the $h(g(p))$, and $h \mapsto hg$ runs over all eight symmetries: the orbit of $q$ is the orbit of $p$. Two orbits that
share an element are therefore both equal to that element's orbit.

\textbf{16.} $g(p) = h(p)$ exactly when $h^{-1}g$ fixes $p$. If $k$ symmetries fix $p$, each image is reached by
exactly $k$ symmetries, and the orbit has $8/k$ elements.

\textbf{17.} The mirror $(m - r, c)$ would need the queens of rows $r$ and $7 - r \neq r$ in the same column. The mirror
$(r, m - c)$ would need $c = 7 - c$, impossible for an integer. Under $(c, r)$, a queen off the main diagonal pairs with
another at the mirror cell, and $r + c = c + r$ puts them on one diagonal; so all eight would sit on the main diagonal,
which holds only one. The other diagonal mirror is the same with $r - c$.

\textbf{18.} The quarter-turn generates four \omterm{def:b1:balanced-trees:rotation}{rotations}; if it fixed $p$, at least four symmetries would fix $p$ and the
orbit would have at most 2 elements. The output has orbits of 8 and of 4 only. The orbit of 4 has a pair of symmetries
fixing it, the identity and one of order 2: not a reflection (question 17), so the half-turn.

\textbf{19.} That the permutation has eight distinct values of $r - c$ and eight distinct values of $r + c$.

\textbf{20.} It moves the eight cells, sorts them by row, and reads off the columns, so the image is written again as
the list of the columns of rows 0 to 7, and equal solutions are equal tuples.

\textbf{21.} It takes the first solution not yet seen, computes the set of its images and marks them all as seen. By
question 15 the orbits are disjoint, so each solution is in exactly one, found when its first member is met.

\textbf{22.} $88 + 4 = 92$.

\textbf{23.} For $(2, 4, 1, 7, 0, 6, 3, 5)$, the half-turn needs $p(7 - r) = 7 - p(r)$: $p(7) = 5 = 7 - 2$, $p(6) = 3 =
7 - 4$, $p(5) = 6 = 7 - 1$, $p(4) = 0 = 7 - 7$, and the other four follow.

\textbf{24.} \emph{The eight-queens count}: eight queens can be placed on a chessboard with no two attacking each other
in 92 ways; up to the eight symmetries of the board there are 12, eleven with eight images and one, fixed by the
half-turn, with four.
\end{solution}
